\author{}
\documentclass[11pt,a4paper]{article}
\usepackage[utf8]{inputenc}
\usepackage[T1]{fontenc}
\usepackage[french]{babel}

% Paquets pour les mathématiques
\usepackage{amsmath, amssymb, array}

% Marges
\usepackage{geometry}
\geometry{hmargin=2cm, vmargin=2cm}

% Paquet pour les boîtes colorées
\usepackage[many]{tcolorbox}

% Définition d'un style personnalisé pour les définitions (bleu)
\newtcolorbox{boiteDef}[1]{
    colback=blue!5!white,      % Couleur de fond (très clair)
    colframe=blue!75!black,    % Couleur du contour et du titre
    title={\textbf{#1}},       % Le titre en gras
    arc=3mm,                   % Arrondi des coins
    boxrule=1pt,               % Épaisseur du trait
    fonttitle=\bfseries        % Police du titre
}

% Définition d'un style pour les propriétés (rose)
\newtcolorbox{boiteProp}[1]{
    colback=pink!5!white,
    colframe=purple!70!pink, 
    title={\textbf{#1}},
    arc=3mm,
    boxrule=1pt,
}

% Définition d'un style pour les rappels (vert)
\newtcolorbox{boiteRappel}[1]{
    colback=green!10!white,
    colframe=green!50!black,  
    title={\textbf{#1}},
    arc=3mm,
    boxrule=1pt,
}

% Définition d'un style pour les remarques (jaune)
\newtcolorbox{boiteRemarque}[1]{
    colback=yellow!10!white,
    colframe=yellow!70!black,  
    title={\textbf{#1}},
    arc=3mm,
    boxrule=1pt,
}

\renewcommand{\thesection}{\Roman{section}}

\begin{document}

\title{\textbf{Chapitre 3 : Limites de suites }}
\date{}
\maketitle

\section{Définitions et premières propriétés}

\begin{boiteDef}{Définition :}
    Une suite $(u_n)$ tend vers +$\infty$ si tout intervalle ouvert de la forme $]A; +\infty[$ avec $A \in \mathbb{R}$ contient tous les termes de la suite à partir d'un certain rang $n_0$.
    \newline
    On note alors $\displaystyle \lim_{n \to +\infty} u_n = +\infty$ ou $u_n \to +\infty$.
\end{boiteDef}

\begin{flushleft}
    En d'autres termes, cela signifie que pour tout nombre réel $A$, il existe un rang $n_0$ tel que pour tout $n \geq n_0$, on a $u_n > A$.
    A partir d'un certain rang, les termes de la suite sont aussi grands que l'on veut.
\end{flushleft}

\begin{boiteProp}{Propriété :}
    Les suites ${\sqrt{n}}$, $(n)$, $(n^p)$ avec $p \in \mathbb{N}^*$ ont pour limite +$\infty$.
\end{boiteProp}

\begin{boiteDef}{Définition :}
    Une suite $(u_n)$ tend vers -$\infty$ si tout intervalle ouvert de la forme $]-\infty; A[$ avec $A \in \mathbb{R}$ contient tous les termes de la suite à partir d'un certain rang $n_0$.
    \newline
    On note alors $\displaystyle \lim_{n \to +\infty} u_n = -\infty$ ou $u_n \to -\infty$.
\end{boiteDef}

\begin{flushleft}
    En d'autres termes, cela signifie que pour tout nombre réel $A$, il existe un rang $n_0$ tel que pour tout $n \geq n_0$, on a $u_n < A$.
    A partir d'un certain rang, les termes de la suite sont aussi petits que l'on veut.
\end{flushleft}

\begin{boiteRemarque}{Remarques :}
    La limite d'une suite se cherche toujours en +$\infty$.
    \newline
    Une suite qui a une limite infinie est dite divergente.
\end{boiteRemarque}

\begin{boiteDef}{Définition :}
    Une suite $(u_n)$ a pour limite le réel $l$ si tout intervalle ouvert contenant $l$ contient tous les termes de la suite $(u_n)$ à partir d'un certain rang.
    \newline
    On note alors $\displaystyle \lim_{n \to +\infty} u_n = l$ ou $u_n \to l$.
    \newline
    On dit aussi que la suite $(u_n)$ converge vers $l$.
\end{boiteDef}

\begin{flushleft}
    En d'autres termes, cela signifie que pour tout nombre réel $\varepsilon > 0$, il existe un rang $n_0$ tel que pour tout $n \geq n_0$, on a $|u_n - l| < \varepsilon$.
    \newline
    A partir d'un certain rang, les termes de la suite sont aussi proches de $l$ que l'on veut.
\end{flushleft}

\begin{boiteProp}{Propriété :}
    Les suites $(\frac{1}{n})$, $(\frac{1}{n^p})$, $(\frac{1}{\sqrt{n}})$ avec $p \in \mathbb{N}^*$ ont pour limite $0$.
\end{boiteProp}

\begin{boiteProp}{Unicité de la limite :}
    Si une suite $(u_n)$ a une limite finie $l$, alors cette limite est unique et on la note .
\end{boiteProp}

\begin{boiteRemarque}{Vocabulaire :}
    Si une suite a une limite finie, on dit que la suite est convergente.
    \newline
    Si une suite n'a pas de limite finie, on dit que la suite est divergente.
\end{boiteRemarque}

\section{Opérations sur les limites}

\subsection{Somme}
\begin{center}
\begin{tabular}{|c|c|c|c|c|c|c|}
\hline
$\displaystyle \lim_{n\to+\infty} u_n$ & $l$ & $l$ & $l$ & $+\infty$ & $-\infty$ & $+\infty$ \\ \hline
$\displaystyle \lim_{n\to+\infty} v_n$ & $l'$ & $+\infty$ & $-\infty$ & $+\infty$ & $-\infty$ & $-\infty$ \\ \hline
$\displaystyle \lim_{n\to+\infty} u_n+v_n$ & $l+l'$ & $+\infty$ & $-\infty$ & $+\infty$ & $-\infty$ & F.I. \\ \hline
\end{tabular}
\end{center}

\subsection{Produit}
\begin{center}
\begin{tabular}{|c|c|c|c|c|c|c|c|}
\hline
$\displaystyle \lim_{n\to+\infty} u_n$ & $l$ & $l>0$ & $l<0$ & $l>0$ & $l<0$ & $\infty$ & $0$ \\ \hline
$\displaystyle \lim_{n\to+\infty} v_n$ & $l'$ & $+\infty$ & $+\infty$ & $-\infty$ & $-\infty$ & $\infty$ & $+\infty$ \\ \hline
$\displaystyle \lim_{n\to+\infty} u_n\times v_n$ & $l\times l'$ & $+\infty$ & $-\infty$ & $-\infty$ & $+\infty$ & $\infty$ signe à définir & F.I. \\ \hline
\end{tabular}
\end{center}

\subsection{Quotient}
\begin{center}
\begin{tabular}{|c|c|c|c|c|c|c|}
\hline
$\displaystyle \lim_{n\to+\infty} u_n$ & $l$ & $l$ & $l\neq 0$ & $0$ & $\infty$ & $\infty$ \\ \hline
$\displaystyle \lim_{n\to+\infty} v_n$ & $l'\neq 0$ & $\infty$ & $0$ & $0$ & $l'\,(\text{évent.})$ & $\infty$ \\ \hline
$\displaystyle \lim_{n\to+\infty} \dfrac{u_n}{v_n}$ & $\dfrac{l}{l'}$ & $0$ & $\infty$ signe à définir & F.I. & $\infty$ signe à définir & F.I. \\ \hline
\end{tabular}
\end{center}

\section{Théorèmes de comparaison}

\subsection{Pour les limites infinies}

\begin{boiteProp}{Théorème}
    Si $(u_n)$ et $(v_n)$ sont deux suites telles que :
    \begin{itemize}
        \item $u_n \leq v_n$ à partir d'un certain rang $n_0$,
        \item $\displaystyle \lim_{n \to +\infty} u_n = +\infty$,
    \end{itemize}
    alors $\displaystyle \lim_{n \to +\infty} v_n = +\infty$.
\end{boiteProp}

\begin{boiteProp}{Corollaire du théorème}
    Si $(u_n)$ et $(v_n)$ sont deux suites telles que :
    \begin{itemize}
        \item $u_n \leq v_n$ à partir d'un certain rang $n_0$,
        \item $\displaystyle \lim_{n \to +\infty} v_n = -\infty$,
    \end{itemize}
    alors $\displaystyle \lim_{n \to +\infty} u_n = -\infty$.
\end{boiteProp}

\subsection{Pour les limites finies}

\begin{boiteProp}{Propriété}
    Si $(u_n)$ a pour limite $l$ et si $(v_n)$ est une suite telle que $u_n \leq v_n$ à partir d'un certain rang $n_0$, alors $\displaystyle \lim_{n \to +\infty} v_n \geq l$.
\end{boiteProp}

\begin{boiteProp}{Cas particuliers}
    Une suite $(u_n)$ à termes positifs ou nuls à partir d'un certain rang a une limite $l \geq 0$.
    \newline
    Une suite $(u_n)$ à termes négatifs ou nuls à partir d'un certain rang a une limite $l \leq 0$.
\end{boiteProp}

\begin{boiteProp}{Théorème des gendarmes}
    Si $(u_n)$, $(v_n)$ et $(w_n)$ sont trois suites telles que :
    \begin{itemize}
        \item $u_n \leq v_n \leq w_n$ à partir d'un certain rang $n_0$,
        \item $\displaystyle \lim_{n \to +\infty} u_n = l$,
        \item $\displaystyle \lim_{n \to +\infty} w_n = l$,
    \end{itemize}
    alors $\displaystyle \lim_{n \to +\infty} v_n = l$.
\end{boiteProp}

\section{Suites géométriques}

\begin{boiteProp}{Propriété}
    Soit $q$ un nombre réel.
    \begin{itemize}
        \item Si $q > 1$, alors la suite $(q^n)$ diverge et sa limite est +$\infty$.
        \item Si $q = 1$, alors la suite $(q^n)$ converge vers $1$.
        \item Si $-1 < q < 1$, alors la suite $(q^n)$ converge vers $0$.
        \item Si $q \leq -1$, alors la suite $(q^n)$ diverge et n'a pas de limite.
    \end{itemize}
\end{boiteProp}

\begin{boiteRemarque}{Remarque :}
    Les suites de type $(q^n)$ sont des suites géométriques de premier terme égal à $1$.
\end{boiteRemarque}

\section{Limite des suites monotones}

\begin{boiteProp}{Propriété}
    Si une suite croissante a pour limite le réel $l$, alors tous les termes de la suite sont majorés par $l$. 
    Si une suite décroissante a pour limite le réel $l$, alors tous les termes de la suite sont minorés par $l$.
\end{boiteProp}

\begin{boiteProp}{Théorème}
    \begin{itemize}
        \item Une suite croissante et majorée converge (limite finie $l \in \mathbb{R}$)
        \item Une suite croissante et non majorée diverge vers +$\infty$.
        \item Une suite décroissante et minorée converge (limite finie $l \in \mathbb{R}$)
        \item Une suite décroissante et non minorée diverge vers -$\infty$.
    \end{itemize}
\end{boiteProp}

\end{document}
